\documentclass[12pt]{article}
\usepackage[utf8]{inputenc}
\usepackage{graphicx}
\usepackage{amsmath}
\usepackage{hyperref}
\usepackage{natbib}
\usepackage{geometry}
\geometry{margin=1in}
\usepackage{subcaption}
\usepackage{pdfpages}

\title{Bike Design Drafts}
\author{Jack Henry}
\date{\today}

\begin{document}



\includepdf[pages=-]{IIA_Project_Coversheet_Feedback_2020}

\maketitle

\section{Introduction}

After initial drafting ideas and a little bit of research, I have grouped the draft plans into 3 main areas: \textbf{automatic transmission}, \textbf{natural materials}, and \textbf{folding bicycles}. 
\\
\\
I think these areas have key opportunities for proof of concepts and innovation.

\section{Automatic Transmission}
Automatic transmission systems in bicycles aim to optimize gear shifting without active input from the rider, improving the ease of use of the bike. These systems can be either electronic or fully mechanical in nature.

\subsection{Mechanical Systems}
One option for implementing an automatic transmission system is using a purely mechanical system. I think there are 2 possible candidates for a mechanical-based system:

\begin{itemize}
    \item \textbf{Centrifugal-based Systems}: Some designs use spinning weights to automatically shift between gears based on wheel speed. \cite{Centrifugal}
    
    \item \textbf{Torque-based Systems}: Some designs rely on torque, using mechanical linkages that expand when more force is gone through them. One interesting concept for this is expandable gears, where the gears get bigger/smaller based on the torque. \cite{Torque} 

\end{itemize}


\begin{figure}[h!]
    \centering

    \begin{subfigure}[b]{0.2\textwidth}
        \centering
        \includegraphics[width=\textwidth]{Centrifugal.png}
        \caption{Centrifugal automatic gears \cite{Centrifugal}}
        \label{fig:centrifugal}
    \end{subfigure}
    \hspace{0.1\textwidth}
    \begin{subfigure}[b]{0.27\textwidth}
        \centering
        \includegraphics[width=\textwidth]{Torque.jpg}
        \caption{Expandable gears concept \cite{Torque}}
        \label{fig:Torque}
    \end{subfigure}
    \hspace{0.1\textwidth}
    \begin{subfigure}[b]{0.5\textwidth}
        \centering
        \includegraphics[width=\textwidth]{Torque2.jpg}
        \caption{Expandable gears concept \cite{Torque}}
        \label{fig:Torque2}
    \end{subfigure}

    \caption{Mechanical automatic gears}
    \label{fig:mechanical}
\end{figure}



These mechanical systems have the advantage of not requiring batteries or electronics, and can generally be low-maintenance and robust. However, if they are not calibrated properly, they can actually make it harder to cycle.

\newpage

\subsection{Electronic Systems}
Modern electronic automatic shifting offers higher precision and more customization. Systems include:

\begin{itemize}
    \item \textbf{Electromechanical Derailleurs}: An electronic based deraileur. Used in systems such as Shimano Di2.\cite{Shimano}
    
    \item \textbf{Sensor-Driven Logic}: Some e-bikes use torque, cadence, and speed sensors to trigger shifting logic based on terrain and rider effort.
\end{itemize}

\subsection{Opportunities for Innovation}
\begin{itemize}
    \item \textbf{Hybrid Mechanical-Electronic Systems}: Combining robust mechanical actuation with sensor-assisted tuning. This can make a reliable mechanical system, with less electronics that are purely used for calibration.
    
    \item \textbf{Low-Cost Mechanical Automation}: Using centrifugal or torque-based gear changers using new materials, 3D printing, new designs etc.
    
    \item \textbf{Software advancements}: Likely not applicable for this project, but there are lots of possible software innovations, eg apps for public bikes, or using AI to anticipate terrain and user behavior and adjust accordingly.
\end{itemize}

\subsection{Challenges}
\begin{itemize}
    \item For mechanical systems: fine-tuning responsiveness and range of adjustment without added complexity.
    \item For electronic systems: ensuring reliability, battery life, and resistance to environmental conditions.
    \item In all cases: balancing cost, weight, and user override/control.
\end{itemize}

\subsection{Current choice for this area}

I would likely opt for exploring expandable gears. I think this is a novel concept that hasn't been explored too much before, and is a somewhat simple solution to a reasonably complex problem. There are many different ways of creating these gears, such as in figure \ref{fig:Linkages} below, and so there would be lots to look into and test.

\begin{figure}[h!]
    \centering

    \begin{subfigure}[b]{0.25\textwidth}
        \centering
        \includegraphics[width=\textwidth]{Linkage.jpeg}
        \caption{Circular linkage \cite{Linkage}}
        \label{fig:Linkage}
    \end{subfigure}
    \hspace{0.05\textwidth}
    \begin{subfigure}[b]{0.3\textwidth}
        \centering
        \includegraphics[width=\textwidth]{Linkage2.png}
        \caption{Expandable linkage \cite{won2022}}
        \label{fig:Linkage2}
    \end{subfigure}
    \hspace{0.05\textwidth}
    \begin{subfigure}[b]{0.3\textwidth}
        \centering
        \includegraphics[width=\textwidth]{Linkage3.jpeg}
        \caption{Cam driven expansion}
        \label{fig:Linkage3}
    \end{subfigure}

    \caption{Examples of expandable linkages that could be used in the automatic transmission mechanism}
    \label{fig:Linkages}
\end{figure}


\section{Natural Materials}
The use of natural or sustainable materials for bicycles is something that isn't new, with the use of wood used previously, however, with increasing awareness about sustainability and the environment, I think exploring a proof of concept for a bicycle built from natural materials would be a valid and interesting project.

\subsection{Current and possible materials to be used}
\begin{itemize}
    \item \textbf{Bamboo}: Bamboo frames offer strength-to-weight efficiency and shock absorption, and are increasingly popular natural material. \cite{BambooBikes}
    \item \textbf{Flax/hemp fiber}: Natural fibre composites are improving and are much more prevalent today than ever. They are eco-friendly replacements for carbon and glass fibre, however are more expensive in the current market. \cite{Flax1}\cite{Flax2}
    \item \textbf{Wooden components}: Using wood in some components can increase sustainability but likely wouldn't be a good area of exploration for this project
    \item \textbf{Naturally derived polymers}: There is increasing research into polymers that are derived from natural sources such as cellulose or fungal-mycelium based polymers. Again, this would likely not be an area of research for this project.
\end{itemize}

\begin{figure}[h!]
    \centering

    \begin{subfigure}[b]{0.3\textwidth}
        \centering
        \includegraphics[width=\textwidth]{Bamboo.jpeg}
        \caption{Bamboo bike \cite{BambooBikes}}
        \label{fig:bamboo}
    \end{subfigure}
    \hspace{0.1\textwidth}
    \begin{subfigure}[b]{0.5\textwidth}
        \centering
        \includegraphics[width=\textwidth]{Flax.jpeg}
        \caption{Flax fibre bike \cite{Flax2}}
        \label{fig:Flax}
    \end{subfigure}

    \caption{Natural materials designs}
    \label{fig:natural}
\end{figure}

\subsection{Research Directions}
\begin{itemize}
    \item \textbf{Natural fibre designs}: Looking into natural fibre designs as a proof of concept and/or exploring designs for ease of manufacture and cost design
    \item \textbf{Combination of materials}: Exploring sandwich composites, mimicking the structure of a snowboard to create a bicycle that has the same performance as a normal metal based one.
\end{itemize}

\subsection{Challenges}
\begin{itemize}
    \item Cost: cost is a large factor for natural fibres, as they are often more expensive than other composites \cite{Flax1}
    \item Design: the absorbing qualities of flax and general behaviour of fibre composites can make flax fibre difficult to work with, paricularly when manufacturing a bicycle.
    \item Structural performance: whilst the performance of natural fibres are improving, they are still weaker than most less sustainable alternatives.
\end{itemize}

\subsection{Current choice for this area}

I would opt for exploring the natural fibre composites. I have more familiarity with them, having made (or attempted to make) a snowboard before using them. There is also proven commercial viability, as shown by companies like Bcomp \cite{Bcomp}, who have build cars from flax fibre.


\begin{figure}[h!]
    \centering

    \begin{subfigure}[b]{0.7\textwidth}
        \centering
        \includegraphics[width=\textwidth]{Car.jpeg}
        \caption{Flax fibre in cars \cite{Bcomp}}
        \label{fig:car}
    \end{subfigure}

    \begin{subfigure}[b]{0.7\textwidth}
        \centering
        \includegraphics[width=\textwidth]{Graph.jpeg}
        \caption{CO2 impact of composite fibres \cite{Bcomp}}
        \label{fig:Graph}
    \end{subfigure}

    \caption{Proven uses and impacts of flax fibres}
    \label{fig:Bcomp}
\end{figure}

\section{Folding Bicycles}
Folding bikes provide compact solutions for commuting and casual use, however can often compromise on performance or ease of use, as well as being expensive.

\subsection{Existing Models}
There are many different existing models. Folding bikes like the Brompton have set benchmarks for portability and design. \cite{Brompton}

\begin{figure}[h!]
    \centering

    \begin{subfigure}[b]{0.4\textwidth}
        \centering
        \includegraphics[width=\textwidth]{Bike.jpeg}
        \caption{Bike mode \cite{Folding}}
        \label{fig:Bikr}
    \end{subfigure}
\hspace{0.1\textwidth}
    \begin{subfigure}[b]{0.25\textwidth}
        \centering
        \includegraphics[width=\textwidth]{Fold.jpeg}
        \caption{Folded \cite{Folding}}
        \label{fig:Fold}
    \end{subfigure}

    \caption{Brompton M6L \cite{Brompton}}
    \label{fig:Bike}
\end{figure}

\subsection{Areas for Development}
There are multiple areas of folding bikes that can be improved, but I highlight 3 below:
\begin{itemize}
    \item \textbf{Wheels}: Wheels are often small, creating instability and other issues. I think an interesting idea is possibly looking into larger wheeled folding bikes - possibly utilizing airless tires to create a tire that can fold itself.
    \item \textbf{Modular, "clipping" bikes}: Rather than a bike that folds into itself, exploring the idea of a bike that can easily be taken apart, and clipped back together. This could improve the robustness of the bike as there would be no hinge points.
    \item \textbf{Backback design}: Designing a folding bike that can fold into a backpack that can easily be carried.
\end{itemize}

\subsection{Challenges}
\begin{itemize}
    \item Ergonomics and ride quality: Folding bikes often have to compromise on these in order to make them sufficiently compact.
    \item Weight: folding bikes are often bought to be carried and transported, so they need to be sufficiently light 
    \item Stability: Folding bikes can often unexpectedly collapse, creating issues such as safety.
\end{itemize}

\subsection{Current choice for this section}
I would opt for exploring the "clip bike" to remove hinges and create a more modular approach to folding bikes. However, I would likely go with one of the the earlier areas.

\section{Other Areas to Explore}
Beyond the three main topics, several other areas may be worth exploring:
\begin{itemize}
    \item \textbf{Smart Bikes}: Embedded software, GPS, anti-theft, and app integration
    \item \textbf{Suspension Design}: Adaptive or electronically controlled damping for hybrid terrain
    \item \textbf{Energy Harvesting}: Use of dynamos or piezoelectric materials to power onboard electronics
    \item \textbf{Self-balancing bike}: Use of gyroscopes and counterweights to create a more stable bike that can stand by itself when at rest.
    \item \textbf{Tandem duality}: A regular bike that can extend to be a tandem
    \item \textbf{Power systems}: Using systems such as pneumatics or flywheels to store energy to be released at time when needed, creating smoother power demands on the rider
    \item \textbf{Snow bike}: A bike that works in the snow, using either phat or boarded wheels.
    

\end{itemize}

\section{Conclusion}

Three main areas of \textbf{automatic transmission}, \textbf{natural materials}, and \textbf{folding mechanisms} were explored. From these, the initial ideas of \textbf{expandable gears}, \textbf{natural fibre bike}, and \textbf{"clip" bike} were developed as the proposed ideas for the project.





\bibliographystyle{plainnat}
\begin{thebibliography}{9}

\bibitem{Centrifugal}
Reddit "Bicycle with an automatic transmission... It uses centrifugal force of the flywheel to change gears." 
\url{https://www.reddit.com/r/bicycling/comments/ht3st1/bicycle_with_an_automatic_transmission_it_uses/}.

\bibitem{Torque}
Arstechnica "Automatic bike transmission concept is wild and spiky—and could be a big shift" \url{https://arstechnica.com/cars/2023/11/automatic-bike-transmission-concept-is-wild-and-spiky-and-could-be-a-big-shift/}.

\bibitem{Shimano}
Shimano "Di2 Electronic Shifting System."  \url{https://bike.shimano.com/technologies/details/di2.html}.

\bibitem{Linkage}
Autodesk "Circular Linkage with Fixed Center Mechanism"  \url{https://www.autodesk.com/community/gallery/project/21713/circular-linkage-with-fixed-center-mechanism-1}.

\bibitem{won2022}
J. Won, S. Ryu, S. Kim, K. Yoo, H. Kim, and T. Seo,
"Design optimization of a linkage-based 2-DOF wheel mechanism for stable step climbing,"
\textit{Scientific Reports}, vol.~12, p.~16912, Oct.~2022.
doi: \href{https://doi.org/10.1038/s41598-022-21410-1}{10.1038/s41598-022-21410-1}.


\bibitem{BambooBikes}
Bamboo Bicycle Club. "Bamboo Bike Kits and Sustainability." \url{https://bamboobicycleclub.org/}.


\bibitem{Flax1}
Easy Composites "Flax Fibre Prices" \url{https://www.easycomposites.co.uk/product-search?keywords=flax}.

\bibitem{Flax2}
Force technology "How to design and produce composite components with natural flax fibres" \url{https://forcetechnology.com/en/articles/learnings-from-designing-and-producing-composite-components-with-natural-fibres}.

\bibitem{Bcomp}
Bcomp \url{https://www.bcomp.com/applications/sports/}.


\bibitem{Brompton}
Brompton Bicycles. "Folding Bike Innovation." \url{https://www.brompton.com/}.

\bibitem{Folding}
The Strategist. "I Tested a Bunch of Folding Bikes to Find the Four Best Ones" \url{https://nymag.com/strategist/article/best-folding-bikes.html}.


\end{thebibliography}


\includepdf[pages=-]{SSD.pdf}

\end{document}
